% SPDX-License-Identifier: CC-BY-4.0
% Copyright (c) 2026 Martin Schröder <info@swedishembedded.com>

\section{Corrections: defects the work exposed, what is fixed and what is open}\label{sec:corrections}

A measured result is never rewritten by a later repair: \cref{sec:results} reports what
the build of the time produced, and this section says what that build got wrong, what was
changed, and what remains. Each repair names the commit that made it; the ledger checks
that the commit is in the history of the pinned revision of its repository (\cref{sec:repro}).
Two sources contribute: the failures of the first registered campaign
(\cref{sec:defects,sec:open}) and a reading of the three code bases against the
requirements of the coding-task protocol before it was fixed (\cref{sec:audit}).

\begin{table}[t]
\centering\small
\begin{tabular}{p{4.7cm}p{3.4cm}p{1.6cm}p{4.2cm}}
\toprule
Defect & Seen in & Status & Commits \\
\midrule
Holdout divided a one-source set & E2/E3 first wave & fixed & 2df4237, 5bef68b, 0e7ffaa, 1aa402b \\
Markup hid the subject & E1 & fixed & 1975575 \\
Adapter card id refused by serve & E2 second wave & fixed & a335ebb \\
Serve bank narrower than training & E2 second wave & fixed & 292512e \\
Thousands separator failed a number & E2 registered run & fixed & 2e72256 \\
Retention blind to taught facts & E2d & open & none \\
Anchor bound too tight for small suites & E2, E2b, E2c & option, default unchanged & cab56de, 6ca8b3f \\
Pass is exit status only & audit & open & none \\
Practice gate runs without canaries & audit & open & none \\
\bottomrule
\end{tabular}
\caption{Defects, where they were seen, and their state at the pinned revisions.}\label{tab:corrections}
\end{table}

\subsection{Defects found in the first campaign and repaired}\label{sec:defects}

Running the loop end to end, rather than reading it, exposed defects between the first
stage and the first release; each was reproduced, pinned by a regression test written
before the repair, and repaired in Splinter. The first wave of E2 stopped at the first,
and the second wave reached the others in turn, because each hid the next.

\paragraph{A set from one source text was divided, then refused.} Records carry the group
of the source text they are grounded in, so a fact sheet holding all facts is one group.
The holdout rule, when no whole group fitted, held out the newest tenth of the records,
dividing the group, and the leakage check on the files about to be written then refused the
rule's own output (``1 group(s) and 0 unit(s) appear in more than one of train, validation
and test''), after every rollout had been paid for. The rule now holds nothing out in that
case and the split names the cause. The experiment design was also at fault: a held-out
measurement needs independent units, so a generation's source is now a directory of one
sheet file per fact.\cl{FX01} Later commits extended the same rule: a free record follows the
fixed side of its group, monitoring with no whole group to set aside is skipped with the
reason recorded, and a source that is one independent text is refused at intake before any
rollout.\cl{FX05}

\paragraph{Markup hid a subject.} The admission rule compared the subject and the instruction
as literal text, so ``the \code{q} key'' did not contain the subject ``q key''. One function
now ignores code quotes, quotation marks, emphasis and a call's empty parentheses for the
instruction, the source and the variants alike.\cl{FX02}

\paragraph{No release could pass the serve check.} The gate's serve check starts
\code{brain serve --adapter}. Training recorded the adapter's card id as
\code{base:splinter:candidate}, one segment short of Brain's adapter grammar, so the server
refused it; after that was repaired it refused the adapter again because Splinter's LoRA
corrects every linear and Brain's adapter bank covers the attention linears unless told
otherwise. The test double of the server had accepted any flags, which is why neither showed
before a real server was started. The card id now ends in a tag and the check passes the
bank option.\cl{FX03,FX04}

\paragraph{A documentation defect.} Checking the release statements against the code
(\cref{sec:gate}) found that the command-line README placed a release under a directory that
no code creates; the manifest is kept in the experience database, the adapter in the
content-addressed artifact store and the only per-alias file is the serving manifest.\cl{X08}
The README was corrected in the change that introduced this paper.

\paragraph{A correct number failed the verifier.} The stated-reference check treated a number
written with a thousands separator (``4,711'') as different from the number (4711), which with
references that carry a unit word the answer omits left tasks untaught before training
(\cref{sec:results}). A number with a separator is now the same number.\cl{FX06}


\subsection{Defects the campaign exposed that are not closed}\label{sec:open}

\paragraph{Retention is blind to taught facts.} In the 1.7B run (E2d) the second candidate
answered \mETwodCandidateGOneGZero{} of the \mETwodAfterGZeroGZero{} first-generation facts the champion answered, on independent probes,
and passed the gate's retention check. The check holds a candidate to each earlier
release's held-out tasks, and for a facts campaign those are facts the release never
taught.\cl{OP01} Workspace families now add their own exams to that suite;\cl{OP01} facts do
not. The result stands as reported (\cref{sec:results}); the defect is open, and until the
retention exam covers demonstrated capabilities a passed retention check on a facts campaign
is not evidence that taught facts were kept.

\paragraph{The anchor bound is too tight for small suites.} The anchor check refuses a
candidate whose accuracy on the frozen suite drops by more than 0.02. On the 40-question
anchor suite a single net lost answer is a drop of 0.025, so the check allows none, and it
refused every 0.6B candidate of the first campaign on that alone (\cref{sec:results}). The
gate can now decide the anchor on answers lost net of answers gained (an option; the default
rate bound is unchanged), its error rates are computed in a test, and a 200-question anchor
suite built by code and a reviewed list exists.\cl{FX07} We did not rerun E2 with either; the
registered outcome, ``no release'', is unchanged.

\paragraph{A pass is an exit status.} The workspace verifier passes a patch only if every
fail-to-pass and pass-to-pass command exits with status 0.\cl{OP02} It does not prove that
the intended tests ran and passed, so a suite that exits zero without running them would
pass. The tampering guard is a separate path check. This is open; the coding-task
experiments of \cref{sec:pending} register conformance measurements for the verifier
(\cref{sec:protocol}) before they count.

\paragraph{The practice gate has no canaries.} The release gate called from \code{practice}
is built without canary tasks, and a gate with an empty canary set skips that check.\cl{OP03}
A gate that cannot grade its canaries is not known to be clean, so a release from
\code{practice} does not yet carry that evidence.

\subsection{Audit of the stack against the coding-task protocol}\label{sec:audit}

Before the protocol of \cref{sec:protocol} was fixed, the three code bases were read against
its requirements. Each finding was located in source; the repair names the commit and a
regression test where the repository has one.

\begin{enumerate}
\item \textbf{Loss masking of tool conversations} (Brain; fixed). Under the Qwen3 template
an empty reasoning block is added to the last assistant turn only, so non-final assistant
turns of a multi-step conversation failed the trainer's boundary check. Tool-calling
conversations now train per assistant turn in both thinking modes.\cl{AU01}
\item \textbf{Tool definitions lost between recording and training} (fixed). The agent
records the tool definitions of the run, but the trajectory projections dropped them; an
ATIF trajectory's tool definitions now become the chat sample's tools, and Splinter's
dataset check covers tools.\cl{AU02}
\item \textbf{Trajectory fidelity in the agent SDK} (Sven; fixed). The exported trajectory
marked no context compaction and kept the uncompacted history, listed every registered tool
rather than those offered, and turned unparseable tool-call arguments into an empty object.
Compactions are now marked and the compacted history kept, the offered tools are exported,
and the raw text of non-JSON arguments is kept.\cl{AU03}
\item \textbf{Release probes were closed-book} (fixed). A coding agent is now gated on the
exam its run reserved, with report types for workspace suites.\cl{AU04}
\item \textbf{Contamination paths} (fixed). The semantic-memory tool is withheld from the
worker, and a context tool that could open any readable path is jailed.\cl{AU05}
\item \textbf{Silent waits} (Sven; fixed). A model server that goes quiet is bounded and
the wait is reported, with a stalling mock provider to test it.\cl{AU06}
\item \textbf{Tool-definition mismatches} (Sven; fixed). Definitions disagreed with the
implementation and were tightened.\cl{AU07}
\item \textbf{Serving several adapters} (Brain; fixed in the qwen3 server). An adapter bank
lets one batch serve the base and many LoRA adapters.\cl{AU08}
\item \textbf{Containment} (open). A coding agent's shell is not a sandbox. Where no
container is configured every environment record says \code{isolation: none}, and the
exposure audit detects exposure but cannot prevent it.\cl{AU09}
\end{enumerate}
